
The overall architecture of OSTRICH2 is depicted schematically in Figure~\ref{fig:architecture}.
In essence, OSTRICH2 comprises three string solvers that can be used independently or run using a time-sliced portfolio approach.

At its core, OSTRICH2 builds on the SMT solver Princess~\cite{princess}, which provides the logical reasoning framework and support for theories such as linear integer arithmetic.
OSTRICH2 provides pre-processing and support for the quantifier-free string theories \verb+QF_S+ (pure strings) and \verb+QF_SLIA+ (strings with linear integer arithmetic).
Princess also provides support for algebraic data-types (ADTs), on top
of which the ADT-Str string solver is implemented.

The second and third solver are the Regular Constraint Propagation
(RCP) and Cost-Enriched String (CE-Str) engine, respectively.
These share similar architectures.
They begin with a shared string preprocessor, which processes the SMT-LIB input and forwards it to the SMT Core of Princess.
Details of the string preprocessing are given in the next section.
The SMT Core coordinates with the LIA Solver for linear integer arithmetic and with a solver engine for string-specific reasoning.

Both solvers uses automata- and tranducer-based representations of sets of strings and supported string functions.
A main loop applies proof rules like Automata Intersection, Forward (FWD) and Backward (BWD) Propagation, Nielsen Splitting, and other string-related inference rules.
Details of the rules are given in the next section.
A string database and an automata database are used for efficiently storing and retrieving string and automata data.
The automata database is based on the BRICS Automata Library~\cite{brics}, which provides efficient automaton operations.

For the purposes of presentation, the solvers can be considered to receive and handle constraints from the SMT Core in the following normal form
\begin{align*}
    S ~~::=~~ & A \;|\; \neg A \;|\; S \land S \\
    A ~~::=~~ &
        f(x_1, \ldots, x_n)
        \;|\; x = g(x_1, \ldots, x_n)
        \;|\; \\ & x \in L
        \;|\; x = \sum_i d_i x_i \ .
\end{align*}
Here $x$, $x_1$, \ldots, $x_n$ are string or integer variables and $d_i$ are integer constants, $f$ denotes a string predicate over parameters $x_1, \ldots, x_n$, and $g$ denotes a string function with input parameters $x_1, \ldots, x_n$ and output $x$.
We will make use of standard string functions and predicates such as $\sconcat$, $\sprefixof$, $\scontains$, and $\ssuffixof$ as needed.
We sometimes write $|t|$ to denote the length of the word represented by the term $t$.
For example, the constraint
\begin{lstlisting}[language=SMTLib,frame=single]
(assert (and (= (str.++ x y) (str.++ y x))
             (str.in_re x (re.from_ecma2020 "a*ba*"))
             (str.in_re y (re.from_ecma2020 "a*ca*"))))
\end{lstlisting}
has the following normal form using an additional variable $z$
\[
    z = \sconcat(x,y) \wedge z = \sconcat(y,x)
        \wedge x \in a^*ba^* \wedge y \in a^*ca^* \ .
\]

%ecma script 2020 -> remove

\begin{figure}
    \begin{center}
        \resizebox{\columnwidth}{!}{ % Scale the diagram to fit within the column width
            \begin{tikzpicture}[node distance=2cm]

                % Nodes
                \node (input) [] {SMT-LIB Input};
                \node (preprocessor) [block, below of=input] {String Preprocessor};
                \node (princess) [below of=preprocessor,yshift=0.2cm] {Princess};
                \node (smtcore) [block,below of=princess, yshift=1.1cm] {SMT Core};
                \node (lia) [block, below of=smtcore, yshift=0.6cm] {LIA/ADT Solver};
                \node (proof) [right of=input, xshift=3cm] {RCP/CEA Rules};
                \node (database) [block, right of=proof, xshift=3cm] {String Database};
                \node (automata) [block, below of=database, yshift=0.6cm] {Automata Database};
                \node (brics) [block, below of=automata] {Brics Automata Library};
                \node (intersection) [block, below of=proof, yshift=0.6cm] {Automata Intersection};
                \node (fwdprop) [block, below of=intersection, yshift=0.6cm] {FWD Propagation};
                \node (bwdprop) [block, below of=fwdprop, yshift=0.6cm] {BWD Propagation};
                \node (nielsen) [block, below of=bwdprop, yshift=0.6cm] {Nielsen Splitting};
                \node (etc)[below of=nielsen, yshift=1.2cm] {etc.};

                % Arrows
                \draw [arrow] (input) -- (preprocessor);
                \draw [arrow, shorten >=0.5cm] (preprocessor) -- (princess);
                \draw [arrow] (smtcore) -- (lia);
                \draw [arrow] (automata) -- (brics);
                \draw [arrow] (brics) -- (automata);

                \draw [arrow, <->] (2,-4) -- (2.8,-4);
                \draw [arrow, shorten >=0.9cm] (database) -- (proof);
                \draw [arrow, shorten <=0.95cm] (proof) -- (database);
                \draw [arrow, shorten >=0.5cm] (automata) -- (intersection);
                \draw [arrow, shorten <=0.55cm] (intersection) -- (automata);
                % Draw the large box around the middle portion
                \node[bigbox, fit=(proof)(intersection)(nielsen)(fwdprop)(bwdprop)] {};
                \node[bigbox, fit=(princess)(smtcore)(lia)] {};

            \end{tikzpicture}
        }
    \end{center}
    \caption{Overall architecture of OSTRICH2: the SMT-LIB input is handled by our string preprocessor and Princess SMT core (with LIA and ADT-Str), while the RCP and CE-Str solvers consist of inference rules that are repeatedly applied using  string and automata databases.}
    \label{fig:architecture}
\end{figure}

In the sections below, we given an overview of the three main solver engines.
These are: ADT-Str (algebraic data-types), RCP (regular constraint propagation), and CE-Str (cost-enriched strings).

\subsection{ADT-Str: List-Based Solver}

The ADT-Str solver builds on the decision procedure for algebraic data-types (ADTs) with size catamorphism implemented in Princess~\cite{DBLP:conf/synasc/HojjatR17}.
Algebraic data-types are used to represent strings using the standard encoding of lists with \texttt{nil} and \texttt{cons} constructors.
The length of a string is computed using the build-in \texttt{size} function provided by the ADT solver, mapping every constructor term to the number constructor occurrences.
Other SMT-LIB functions on strings, for instance substring, concatenation, etc., are in ADT-Str encoded using uninterpreted functions and axioms capturing the recursive
definition of the string functions.
Regular expression matching is implemented using Brzozowski derivatives~\cite{Brzozowski}.

ADT-Str is  useful, in particular, for computing solutions of string
constraints that are outside of the fragments for which the other
solvers are complete.
ADT-Str can easily handle certain functions that are hard for our propagation algorithms.
Those functions include, among others, string-to-integer conversion, and functions like substring and indexof that calculate with integer offsets.

\subsection{RCP: Regular Constraint Propagation}

\emph{Regular Constraint Propagation (RCP)} is the newest algorithm that has been implemented in OSTRICH2, based on a subset of proof rules in our paper~\cite{popl22}.
The main goal of RCP is to \emph{prove unsatisfiability of the input constraint}.
The algorithm handles string functions like concatenation, replace, replaceall, and regular constraints.
Other string functions (e.g.\ reverse, one-way and two-way transducers) are also permitted.
These functions permit either exact pre-image or exact post-image computation of regular constraints or both.
In general, these images need not be exact, but must at least overapproximate the true pre/post image.

The main idea behind RCP is to propagate regular constraints of the form
$x \in L$ to other string variables through forward and backward propagations using the RCP inference rule described in the next section.

\subsection{CE-Str: Cost-Enriched String Engine}

The main algorithm of CE-Str (based on \cite{atva2020}) applies backward regular constraint propagation.
It uses cost-enriched finite automata (CEFAs) instead of standard finite automata to represent sets of words.
These automata are able to capture numerical information about accepted strings, such as the string length, the number of characters appearing before a transition is fired, and so on.
This allows length constraints on strings and functions like indexof to be directly supported.

In addition, CE-Str can handle counting operators more efficiently by leveraging the numerical information encoded in CEFAs. For instance, the regex $a^{\{1,1000\}}$, which accepts strings consisting of 1 to 1000 repetitions of the character $a$, can be represented as a CEFA with just one state and one transition—compared to the 1000 states and 1000 transitions required by a standard finite automaton.

CE-Str supports a wide range of string functions, including concatenation, replace, replaceall, substring, indexof and length. It provides completeness guarantees for the straight-line fragment of these functions. CE-Str complements RCP and ADT-Str. It performs well on string constraints involving integer type but struggles when the constraints are outside its decidable fragment.

